% Volume IX: Institutional Design
% Part XVII. Designing for Delayed Truth
% Chapter 103: Audit Without Total Surveillance
% Spine: proof-spines/ch103-spine.md

\chapter{Audit Without Total Surveillance}
\label{ch:ch103}

Accountability does not always require universal exposure of the personal data on which institutions act. A system can be audited through selective audit, cryptographic commitments, access logs, independent custodians, zero-knowledge methods, statistical inspection, and trigger-based disclosure. These mechanisms make it possible to test whether officials acted lawfully without making every subject fully visible.

The design aim is to separate oversight from indiscriminate surveillance. Audit should track who accessed what, when, under what authority, and with what downstream effect, while keeping ordinary persons from becoming permanently inspectable merely because institutions need to be checked. The result is a form of answerability that does not reproduce the panopticon in the name of control.
